\section{Methods}
\label{sec:methods}

\subsection{Common Design Elements}
\label{sec:methods-common}

Three commercially deployed LLMs were evaluated across all studies: Claude
(claude-sonnet-4-20250514; Anthropic), GPT-4 (gpt-4o; OpenAI), and Gemini
(gemini-2.5-flash; Google DeepMind). All models were queried via their public APIs
at default temperature with a maximum of 512 output tokens per response.

Each prompt provided the current value of the target quantity and requested a point
estimate plus 50\%, 80\%, and 90\% confidence intervals as structured JSON output,
with a specified target horizon. No chain-of-thought or explicit calibration
instructions were given; models were evaluated in their default,
deployment-ready configuration. Three confidence levels (50\%, 80\%, 90\%) were
used in all studies. Calibration was measured via Expected Calibration Error (ECE),
defined as the mean absolute deviation between stated and empirical coverage
across all three levels (see Section~\ref{sec:methods-metrics}).

\subsection{Study 1: Equities, Single Horizon}
\label{sec:methods-study1}

Study 1 (March 23, 2026) employed 20 large-cap S\&P 500 equities (AAPL, NVDA,
MSFT, AMZN, GOOGL, META, AVGO, TSLA, BRK-B, WMT, LLY, JPM, XOM, V, JNJ, MU,
MA, ORCL, COST, SPY) with a 1-day prediction horizon (outcomes collected March
24, 2026). Each model made 5 independent predictions per ticker, yielding 100
scored predictions per model. This study establishes a baseline calibration
estimate under the simplest possible conditions: a single, well-defined domain
and the shortest practical horizon.

\subsection{Study 2: Equities, Multi-Horizon}
\label{sec:methods-study2}

Study 2 expanded to 100 large-cap equities across nine prediction horizons (1,
2, 3, 6, 7, 18, 20, 21, and 22 days), with a reference date of March 24, 2026
and predictions made against historical prices via backtesting. Five independent
runs per ticker per model were collected at each horizon. After excluding
unresolved predictions due to market closures (applied tolerance window of 6
calendar days), the study yielded between 157 and 1,497 scored predictions per
model-horizon cell (total $N > 18{,}000$). This design allows direct measurement
of how confidence interval miscalibration scales with the prediction horizon.

\subsection{Study 3: Multi-Domain, Single Horizon}
\label{sec:methods-study3}

Study 3 (March 25, 2026) extended prediction to six heterogeneous domains:
equities ($n = 20$ large-cap tickers), commodity futures ($n = 5$: gold, crude oil,
silver, natural gas, copper), cryptocurrency ($n = 10$ by market cap), foreign
exchange ($n = 8$ major pairs), weather ($n = 30$ U.S.\ cities, next-day high
temperature), and NBA game totals ($n = 3$). All predictions used a 24-hour
horizon. Outcomes were collected automatically March 26, 2026. CoinGecko rate
limits caused five cryptocurrency outcomes to remain unresolved; NBA games did
not achieve completed status in the ESPN API and were excluded. The final scored
dataset comprised $n = 68$ resolved predictions per model ($N = 204$ total).

\subsection{Calibration Metrics}
\label{sec:methods-metrics}

\textbf{Hit rate.} For each confidence level $\alpha \in \{0.50, 0.80, 0.90\}$,
the empirical hit rate is the proportion of predictions for which the true outcome
fell within the stated interval $[L, U]$. A calibrated model achieves a hit rate
equal to the stated confidence; values below indicate overconfidence, values above
indicate underconfidence.

\textbf{Expected Calibration Error (ECE)} is the mean absolute deviation between
stated and empirical coverage across all three confidence levels:
\begin{equation}
  \label{eq:ece}
  \text{ECE} = \frac{1}{3} \sum_{\alpha} \left| \alpha - \hat{p}_\alpha \right|,
  \quad \alpha \in \{0.50, 0.80, 0.90\}.
\end{equation}

\textbf{Meta-knowledge ratio (Study 3 only).} We apply the MEAD/MAD framework of
\citet{Soll2004}. The 80\% CI endpoints are treated as the 10th and 90th percentiles
of a symmetric normal ($z = 1.2816$). The Absolute Deviation (AD) is the midpoint
error of the 80\% CI relative to the true outcome, standardized by within-domain
outcome standard deviation $\sigma$:
\begin{equation}
  \label{eq:ad}
  \text{AD}_i = \frac{|(L_i + U_i)/2 - y_i|}{\sigma}.
\end{equation}

The Expected Absolute Deviation (EAD) is the CI-width-implied error standardized
by $\sigma$, and the meta-knowledge ratio $\mu = \text{MEAD}/\text{MAD}$ is the
ratio of their domain means. $\mu < 1$ indicates overconfidence; $\mu > 1$
indicates underconfidence.
